% =============================================================================
% Module 20: Solutions to Exercises
% =============================================================================
% This file is conditionally included from main.tex via \ifsolutions
% Attested by Solutions/20_Cosmological_Sensitivity/problems_20.wl
% =============================================================================

\section{Solutions to Exercises}
\label{sec:cs_solutions}

\noindent
Full worked solutions are also available in the Mathematica script
\solnlink{20_Cosmological_Sensitivity/problems\_20.wl}{problems\_20.wl},
which verifies each answer symbolically and numerically. The script
computes all distances itself from the flat CPL Friedmann equation, with
an implementation independent of the chapter's companion script. At the
illustrative geometry $(z_d,z_1,z_2)=(0.5,1,2)$ and
$(\Omegam,w_0,w_a)=(0.3,-1,0)$ it finds
$J=(S_{\Omegam},S_{w_0},S_{w_a})=(-0.03569,-0.03301,-0.007255)$, in
agreement with eq.~\eqref{eq:cs_J_values}.
\emph{[Attested: \solnlink{20_Cosmological_Sensitivity/problems\_20.wl}{problems\_20.wl}, check 0a.]}

% ----- Exercise 20.1 -----
\subsection*{Exercise~\ref{ex:cs_cpl}: The CPL Density}

\textbf{(a)} With $\ln a=-\ln(1+z)$, $\dd\ln a=-\dd z/(1+z)$ and the
fluid equation becomes $\dd\ln\rho/\dd z=3[1+w(z)]/(1+z)$, with
$w(z)=w_0+w_az/(1+z)$. Integrating,
\begin{align*}
  \ln\zeta(z)
  &=3(1+w_0)\int_0^z\frac{\dd z'}{1+z'}
   +3w_a\int_0^z\frac{z'\,\dd z'}{(1+z')^2}\\
  &=3(1+w_0)\ln(1+z)+3w_a\Bigl[\ln(1+z)-\frac{z}{1+z}\Bigr],
\end{align*}
which exponentiates to eq.~\eqref{eq:cs_zeta}.

\textbf{(b)} As $z\to\infty$, $z/(1+z)\to1$, so
$\exp[-3w_az/(1+z)]\to e^{-3w_a}$ and
$\zeta\to e^{-3w_a}(1+z)^{3(1+w_0+w_a)}$. The exponent is that of a
constant-$w$ fluid with $w=w_0+w_a$, the high-redshift limit of
eq.~\eqref{eq:cs_cpl_w}.

\textbf{(c)} $\zeta(1)=2^{3(1-0.9-0.3)}\,e^{-3(-0.3)/2}
=2^{-0.6}\,e^{0.45}=1.0347$. The dark-energy density at $z=1$ is only
$3.5\%$ above its present value: $w(z)$ crosses $-1$ at $z=0.5$, so the
density rises going back from $z=0$ to $0.5$ (where $w>-1$) and falls
again beyond (where $w<-1$).

\textbf{(d)} $w_a$ enters only through $w_az/(1+z)$, and
$z/(1+z)\le1/2$ for $z\le1$ and $\le2/3$ for $z\le2$. Over the
redshifts that set $\chi_d$ and $\chi_1$, a unit change in $w_a$ is
therefore equivalent to a much smaller change in $w_0$. At the
illustrative geometry $S_{w_a}/S_{w_0}=0.22$.
\emph{[Attested: \solnlink{20_Cosmological_Sensitivity/problems\_20.wl}{problems\_20.wl}, checks 1a--1f.]}

% ----- Exercise 20.2 -----
\subsection*{Exercise~\ref{ex:cs_projection}: Projection}

For one plane $J=(S_{\Omegam},S_{w_0})$ and $J\cdot\Delta\bm{p}=\delta$
is one equation in two unknowns. For any $\hat{\bm{u}}$ with
$J\cdot\hat{\bm{u}}\neq0$, $\Delta\bm{p}=\delta\hat{\bm{u}}/(J\cdot
\hat{\bm{u}})$ satisfies $J\cdot\Delta\bm{p}=\delta$ identically. With
$\hat{\bm{u}}=(0,1)$, $J\cdot\hat{\bm{u}}=S_{w_0}$ and
$\Delta\bm{p}=(0,\delta/S_{w_0})$. With $\hat{\bm{u}}\propto(c,1)$,
\[
  \Delta w_0=\frac{\delta}{S_{w_0}+cS_{\Omegam}} .
\]
At $\delta=1\%$: $|\Delta w_0|=0.303$ at $c=0$, $0.197$ at
$c=+\tfrac12$ and $0.659$ at $c=-\tfrac12$. Both sensitivities are
negative, so a positive $c$ enlarges the denominator and shrinks the
bias, while a negative $c$ partially cancels $S_{w_0}$ and inflates it.

The fixed-$\Omegam$ value is a lower bound only when
$|S_{w_0}+cS_{\Omegam}|\le|S_{w_0}|$, i.e.\ for $c$ between $0$ and
$2c_\star=-2S_{w_0}/S_{\Omegam}=-1.850$. At $c=2c_\star$ the two are
equal in magnitude, and for $c<2c_\star$ or $c>0$ the marginalised bias
is smaller. ``Marginalising always makes it worse'' is true of
statistical variances (the variance with a parameter fixed never exceeds
the marginalised one), not of biases: a bias moves along
whatever direction the other constraints leave open, and that direction
can either amplify or suppress its component along $w_0$. The
physically relevant comparison is with the degeneracy direction of the
probe combination actually used.
\emph{[Attested: \solnlink{20_Cosmological_Sensitivity/problems\_20.wl}{problems\_20.wl}, checks 2a--2e.]}

% ----- Exercise 20.3 -----
\subsection*{Exercise~\ref{ex:cs_rank}: Rank}

With one plane, $C^{-1}=1/\sigma^2$ and
$J^{\mathsf T}C^{-1}J=\sigma^{-2}J^{\mathsf T}J$ is an outer product
of one vector with itself. Its column space is spanned by $J^{\mathsf T}$
alone, so its rank is $1$ whether $J$ has two or three components. The
null direction in the $(\Omegam,w_0)$ plane is
$\bm{n}\propto(S_{w_0},-S_{\Omegam})=(-0.03301,+0.03569)$, at
$132.8^\circ$ from the $+\Omegam$ axis; along it
$\dd w_0/\dd\Omegam=-S_{\Omegam}/S_{w_0}=-1.081$.

A time delay in the same system raises the rank. It measures
$D_{\Delta t}=(1+z_d)\Dd\Ds/\Dds$ (eq.~\ref{eq:time_delay_recap});
the constant prefactor $(1+z_d)$ drops out of the log-derivatives.
For a single source at $z_s=1$ behind the $z_d=0.5$ deflector, and at
fixed $\Hubble$, its row is
$\partial\ln D_{\Delta t}/\partial(\Omegam,w_0)=(-0.0174,-0.1043)$,
which makes an angle of $38^\circ$ with the DSPL row. The two rows are
independent, so together they have rank~$2$. (The time delay also
depends on $\Hubble$, which then has to be fitted too.) An external
probe in the same likelihood, or a second lens at a different $z_d$,
would also add an independent row; more images of the same two source
planes would not.
\emph{[Attested: \solnlink{20_Cosmological_Sensitivity/problems\_20.wl}{problems\_20.wl}, checks 3a--3c.]}

% ----- Exercise 20.4 -----
\subsection*{Exercise~\ref{ex:cs_coherence}: Coherence}

At $\gp=1.6$ the exponent is $(3-1.6)/(1.6-1)=7/3$, so
\[
  \frac{\mBH_2}{\mBH_1}=0.7^{7/3}=0.43507,\qquad
  \frac{\mBH_3}{\mBH_1}=0.5^{7/3}=0.19843 .
\]
Both are positive and $0<\mBH_3/\mBH_1<\mBH_2/\mBH_1<1$: the three
fractions share the sign of the underlying central mass, and farther
planes feel a weaker version of the same perturbation. They are one
physical error seen at three lever arms, not three independent draws.

By eq.~\eqref{eq:cs_convex} the combined bias in a single parameter is a
weighted mean, with positive weights, of the per-plane biases $b_i$. A
weighted mean of same-signed numbers lies between the smallest and the
largest of them, so the planes cannot cancel each other: the combined
bias is at least as large as the smallest single-plane bias. In the
multi-parameter case the outcome depends on how $\bm{r}$ aligns with the
eigendirections of $F$ (eq.~\ref{eq:cs_coherentbias}), but the
$1/\sqrt{N}$ reduction of independent errors is never available.
\emph{[Attested: \solnlink{20_Cosmological_Sensitivity/problems\_20.wl}{problems\_20.wl}, checks 4a--4d.]}

% ----- Exercise 20.5 -----
\subsection*{Exercise~\ref{ex:cs_nonlinear}: Non-linearity}

The arithmetic gives $0.25/|S_{w_0}|=7.57$ at the illustrative geometry,
and the number is meaningless for two separate reasons.

First, the number is not a displacement. With the radial image masked, the data
do not determine $\gp$, and through $\Brat=\etamea^{\gp-1}$ they do not
determine $\Brat$ either. The $25\%$ is the span $\Brat$ may occupy, and
where the analysis lands inside it is set by the $\gp$ prior. There is
no bias to propagate, because there is no measurement of $\Brat$ to be
biased.

Second, it is off the model manifold. Over the box
$\Omegam\in[0.2,0.4]$, $w_0\in[-2.5,-0.5]$ ($w_a=0$), $\Brat$ varies by
only $[-1.7\%,+3.3\%]$ at this geometry, about eight times less than
$25\%$; and along $\Omegam=0.3$ it is not even monotonic in $w_0$. A
first-order coefficient evaluated at $w_0=-1$ cannot be extrapolated
across such an excursion.

What should be reported is the span of $\Brat$ together with the prior
that selects within it, and the statement that the system's cosmological
constraint has been given up rather than shifted. If a parameter-space
statement is wanted, solve $\Brat(\bm{p})=\Brat_{\rm obs}$ non-linearly
over the span and show the resulting region. If part of the span lies
outside every model prediction, say so: that is a statement that no
cosmology in the family fits.
\emph{[Attested: \solnlink{20_Cosmological_Sensitivity/problems\_20.wl}{problems\_20.wl}, checks 5a--5b.]}

% ----- Exercise 20.6 -----
\subsection*{Exercise~\ref{ex:cs_design}: Design}

The better choice is one well-modelled system, for the following
reasons. Three masked systems add three rows to $J$, so on paper
the rank problem improves. But each masked system has lost the datum that
fixes $\gp$, so each contributes a span of $\Brat$ rather than a value.
Three spans, each selected by a $\gp$ prior, do not converge on the
truth; they converge on whatever the priors agree on.

Nor are the three independent: if they are analysed with the same rigid
profile and the same $\gp$ prior, their errors are shared. The
$\sqrt N$ intuition assumes independent draws; a shared modelling choice
is a shared systematic, and the mean of $N$ identical biases is that
same bias for every $N$.

The single well-modelled system delivers an actual measurement,
with a residual bias that can be quantified and propagated with the
machinery of this chapter. The design lesson is to invest in resolving
and modelling the radial image, rather than in collecting more systems
whose radial images will be masked.
\emph{[Attested: \solnlink{20_Cosmological_Sensitivity/problems\_20.wl}{problems\_20.wl}, check 6a (shared-bias mean); the rest is conceptual.]}

% ----- Exercise 20.7 -----
\subsection*{Exercise~\ref{ex:cs_newgeom}: A Second Geometry}

With the same code as for the illustrative geometry, at
$(z_d,z_1,z_2)=(0.3,0.8,3.0)$:
\[
  \Brat=0.7038,\qquad
  J=(-0.03960,\;-0.04086,\;-0.007386),\qquad
  c_\star=-1.032 .
\]
Compared with $(0.5,1,2)$, the sensitivity to $w_0$ is about $24\%$
larger, and the blind direction has rotated from $c_\star=-0.925$ to
$-1.032$. Neither the size of $S_w$ nor the orientation of the
degeneracy is a property of DSPLs in general; both belong to the
redshift configuration. A conversion factor such as $1/S_w$, or an
amplification factor $\mathcal{A}(c)$, must be recomputed for every
system and never borrowed from another.
\emph{[Attested: \solnlink{20_Cosmological_Sensitivity/problems\_20.wl}{problems\_20.wl}, checks 7a--7d.]}
